\documentclass[12pt]{article}
\usepackage[utf8]{inputenc}
\usepackage[T1]{fontenc}
\usepackage{hyperref}
\title{Patwin Language Resources for K-12 Curriculum Development}
\author{Compiled by Claude Sonnet 5.5 (claude-sonnet-5-5) for Andr\'as Kornai}
\date{}
\begin{document}
\maketitle

\section{Introduction}
Patwin (ISO 639-3 pwi) is a Wintuan language of northern California. By 1997 only one fluent speaker remained at Yocha Dehe, and Lewis Lawyer reports one living first-language speaker today, along with many learners and teachers across three federally recognized tribes. This dossier therefore concerns a heritage and revitalization curriculum.

\section{Primary Educational Resources}
No K-12 textbooks were found in search results. The tribe's own materials (dictionary, grammar book, word database) may be adaptable for teaching, subject to the tribe's permission.

\section{Linguistic and Reference Materials}
\begin{itemize}
\item Lewis C. Lawyer, \emph{A Grammar of Patwin} (University of Nebraska Press, 2021), the first published description of the language. \url{https://muse.jhu.edu/book/82701}
\item UC Davis Magazine, ``Linguistics Alum Writes First Grammar of Patwin''. \url{https://www.ucdavis.edu/magazine/linguistics-alum-writes-first-grammar-of-patwin}
\item Lewis C. Lawyer, UC Davis profile; mentions his Patwin dictionary, privately published by the Yocha Dehe Wintun Nation (year given as 2021 or 2022), so it may not be publicly available. \url{https://nas.ucdavis.edu/people/lewis-lawyer}
\item NSF grant record BCS-1264305 (documentation project). \url{https://grantome.com/grant/NSF/BCS-1264305}
\end{itemize}

\section{Religious and Community Resources}
\begin{itemize}
\item Yocha Dehe Wintun Nation, ``Language is the Heart of Our Culture'': describes the tribe's language programme since 2012, including a dictionary, a grammar book, a database of over 8,000 entries and a tribal language website. \url{https://yochadehe.gov/culture/language/}
\end{itemize}

\section{Digital and Audio Resources}
Yocha Dehe reports an audio and word database and a language website for tribal citizens (see above); access appears to be limited to tribal citizens, and no public audio was seen in search results.

\section{Curriculum Development Context}
Patwin is dormant as a community language but has active revitalization work led by the Yocha Dehe Wintun Nation and other Patwin tribes. Any school curriculum should be developed in partnership with these tribes, which control access to their dictionary and recordings.
\end{document}